% Zdroj: PDF priklady-jaro-2025.pdf, fyzická str. 2-3. Značka: společné okruhy I3. Zařazeno: I3.
\section{Knihovna pro výpočty s maticemi}
\szzotazka{jaro 2025}{2}{Knihovna pro výpočty s maticemi (kontrakt/rozhraní)}

\noindent\textbf{Zadání.}\par
\textit{Poznámky, jak odpovídat na otázku:}
\begin{itemize}
\item Pro odpověď si vyberte jeden z jazyků Java, C++ nebo C\#.
\item Snažte se použít co možná nejnovější syntaxi (a idiomy, atd.) vybraného jazyka.
\end{itemize}

Představte si, že vyvíjíte \textbf{knihovnu pro výpočty s maticemi}.

\begin{enumerate}
\item Definujte kontrakt (bez implementace) pro reprezentaci dvourozměrné matice. Tento
  kontrakt musí poskytovat \textit{prostředky} (tj. metody, funkce, vlastnosti apod.
  podle zvoleného jazyka) alespoň pro:
  \begin{itemize}
  \item získání velikosti jednotlivých rozměrů,
  \item získání hodnoty z matice na základě souřadnic dané hodnoty,
  \item provádění aritmetických operací (sčítání s jinou maticí, odčítání, násobení, atd.).
  \end{itemize}

\item Vytvořte implementaci svého kontraktu, která je vhodná pro velmi velké, ale
  \textbf{řídké} (tzv. \textit{sparse}) matice (tj. matice obsahující převážně nuly).
  Okomentujte svou implementaci tak, aby bylo zřejmé, jak funguje. Z metod (funkcí, atd.)
  implementujte pouze ty, které vrací velikosti rozměrů a hodnotu z matice na základě
  souřadnic.

  Definujte prostorovou složitost své implementace ve stylu $O(\textit{something})$.

\item Představte si, že máte více implementací svého kontraktu pro matice (např. pro řídké
  matice a pro husté matice). Dále si představte, že máte metodu (nebo funkci, atd.),
  která dostane soubor s hodnotami matice a určí, která implementace je pro danou matici
  ideální.

  S využitím vhodného \textit{návrhového vzoru} implementujte \textit{prostředek}, který
  dostane soubor a vrátí instanci kontraktu matice s ideální implementací.
\end{enumerate}

\medskip
\noindent\textbf{Náčrt řešení.}\par
\begin{enumerate}
\item Rozhraní nebo abstraktní třída s metodami podle otázky, přetíženými operátory
  (pokud jsou podporovány), vlastnostmi (pokud jsou podporovány) atd.

\item Existuje mnoho možných implementací (pro matici M$\times$N určitě není správnou
  implementací dvourozměrné pole M$\times$N, které obsahuje všechny hodnoty). Správné
  řešení je uložit pouze nenulové hodnoty (plus nějaká „management" data). Možná řešení:
  (i)~seznam obsahující trojice [řádek, sloupec, hodnota] se seřazením podle řádkového
  indexu a pak podle sloupcového indexu, (ii)~seznam seznamů obsahující dvojice [sloupec,
  hodnota] a seřazený, (iii)~Compressed Sparse Row formát (CSR), tj. tři pole, z nichž
  jedno drží nenulové hodnoty, druhé drží sloupec, kde se daná hodnota nachází, a třetí
  drží index (v hodnotách a sloupcích), kde daný řádek začíná, (iv)~atd. Implementace
  metody getValue závisí na zvolené reprezentaci. Pro CSR to může být například:

\begin{lstlisting}[language=Java]
public double getValue(int row, int col) {
    checkBounds(row, col);
    int start = rowPointers[row];
    int end = rowPointers[row + 1];
    for (int idx = start; idx < end; idx++) {
        if (colIndices[idx] == col) {
            return values[idx];
        }
    }
    return 0.0;
}
\end{lstlisting}

  {\sloppy Složitost závisí na zvolené implementaci, ale měla by být podobná
  $O(\textit{pocetNenulovychHodnot} + \textit{neco})$.\par}

\item Odkazuje to na implementaci na základě návrhového vzoru \textit{Factory}, například:

\begin{lstlisting}[language=Java]
public static Matrix create(Path path) {
    return switch (getSuitableType(path)) {
        case "sparse": new SparseMatrix(path);
        case "dense": new DenseMatrix(path);
        default throw new RuntimeException("unknown_implementation");
    }
}
\end{lstlisting}
\end{enumerate}

\noindent\textit{Doplňující vysvětlení (nad rámec oficiálního náčrtu):} úryvek v~oficiálním náčrtu
míchá dvě syntaxe \texttt{switch} a~v~Javě by se nepřeložil. Výraz \texttt{switch} (Java~14+) se
šipkovými větvemi vypadá takto:
\begin{lstlisting}[language=Java]
return switch (getSuitableType(path)) {
    case "sparse" -> new SparseMatrix(path);
    case "dense"  -> new DenseMatrix(path);
    default -> throw new RuntimeException("unknown implementation");
};
\end{lstlisting}
Jde o~\emph{tovární metodu} (Factory): volající dostane instanci kontraktu \texttt{Matrix}
a~o~konkrétní implementaci (řídká/hustá) rozhoduje jediné místo, které lze rozšířit o~další
reprezentace, aniž se změní kód volajících. Prostorová složitost CSR je $O(\mathit{nnz} + M)$
($\mathit{nnz}$ nenulových hodnot a~sloupcových indexů, $M+1$ ukazatelů na začátky řádků).

\medskip
\begin{csalt}[rozhraní matice + přetížení operátorů]
Více reprezentací (hustá, řídká/CSR) za jedním rozhraním je v~C\# čisté; navíc C\# (na rozdíl
od Javy) umí \emph{přetížení operátorů}, takže \texttt{a + b} funguje jako v~C++:
\begin{lstlisting}[style=cs]
interface IMatrix {
    int Rows { get; } int Cols { get; }
    double this[int r, int c] { get; }            // indexer misto getValue(r,c)
}
sealed class CsrMatrix : IMatrix {
    public int Rows { get; } public int Cols { get; }
    public double this[int r, int c] { get => /* prohledej radek r v colIndices */ 0.0; }
    public static CsrMatrix operator +(CsrMatrix a, CsrMatrix b) => /* ... */ a;
}
static IMatrix Create(string path) => GetSuitableType(path) switch {   // tovarni metoda
    "sparse" => new CsrMatrix(path),
    "dense"  => new DenseMatrix(path),
    _ => throw new ArgumentException("unknown format"),
};
\end{lstlisting}
\end{csalt}
